\noindent{\small\bfseries\textit{Abstract}---GhostVault is a desktop application that moves files
off a computer into private GitHub repositories that the user owns, and brings them back on
request. A file is split into 512\,KiB chunks, and each chunk is encrypted on the computer with
AES-256-GCM under a random key used for that file alone; its associated data binds the chunk to
the file, its position, the number of chunks, the file's exact size and whether it is padded. The
chunks are stored as Git blobs, with one commit per repository. The original is deleted only after
each repository's branch has been read back and lists every chunk under the content id computed
before the upload, a sample of chunks has been downloaded and compared, and the file on disk has
been shown, once more, to be the one that was uploaded. Every step that changes something outside
the application is preceded by a durable journal entry, so an interruption at any point is
finished or undone at the next start. Encrypted recovery files and an optional encrypted backup
kept in the repositories protect against losing the computer. This report describes the system:
its parts and formats, its operations and their failure handling, and where its protection ends,
separating what the cryptography guarantees from what depends on GitHub's access control and on
GitHub keeping the data. It also reports how the system was checked: an automated suite of 663
tests, a randomised run of 3\,200 operations with 1\,273 simulated crashes and 7\,736 injected
network faults in which no file was lost, and measurements of its own costs.\par}

\vspace{6pt}
\noindent{\small\bfseries\textit{Index Terms}---Encrypted storage, authenticated encryption,
AES-GCM, Git object storage, GitHub, crash consistency, write-ahead journaling, desktop
application.\par}
\vspace{4pt}
